\documentclass[12pt,a4paper]{article}
\usepackage[margin=1.5cm]{geometry}
\usepackage[T1]{fontenc}
\usepackage[utf8]{inputenc}
\usepackage{lmodern}
\usepackage{booktabs}
\usepackage{tabularx}
\usepackage{array}
\usepackage{enumitem}
\usepackage{xcolor}
\usepackage{hyperref}
\usepackage{microtype}

\definecolor{brand}{HTML}{176B52}
\definecolor{muted}{HTML}{5F716B}
\hypersetup{colorlinks=true,linkcolor=brand,urlcolor=brand,citecolor=brand}
\setlength{\parindent}{0pt}
\setlength{\parskip}{0.55em}
\newcommand{\status}[1]{\textbf{\color{brand}[#1]}}
\newcommand{\code}[1]{\texttt{#1}}

\title{\textbf{Footprint Lens: Contextual Choice Coach}\\[0.4em]
  \large COMP2300 Deliverable 2: Conceptual Design Report\\
  \normalsize India-Centric Carbon Footprint Awareness and Actionable Reduction}
\author{Arnav Mishra \and Keshav Singhal \and Hiranya Makkar\\[0.4em]
  \small Department of Electronics and Computer Science, University of Southampton}
\date{November 2026}

\begin{document}
\maketitle

\begin{abstract}
\noindent Deliverable 2 presents the conceptual design for \textbf{Footprint Lens}, an interactive system developed for COMP2300 Topic 1. Grounded in our Deliverable 1 empirical findings, the design addresses the core challenge of carbon footprint awareness and reduction within urban Indian households. We report a divergent ideation process yielding eight response families, synthesized into three distinct conceptual models: the \emph{Contextual Choice Coach} (Concept A), the \emph{Household Commons Board} (Concept B), and the \emph{Quiet-Week Action Cards} (Concept C). Through a structured multi-criteria decision analysis evaluating stakeholder fit, feasibility, and risk boundaries, Concept A was selected for prototyping. We present a formal set of small user stories structured across epics, personas, and MoSCoW priorities with testable acceptance criteria. Finally, we document an editable, accessible, low-fidelity prototype spanning seven key interaction states (P0--P6), prove its technical feasibility without speculative machine learning, establish its novelty against existing market tools, and outline a rigorous usability evaluation plan.
\end{abstract}

\hrule
\vspace{0.8em}

\section{Introduction, Problem Framing and Design Principles}

\subsection{Context and Problem Statement}
The COMP2300 Topic 1 brief requires an interactive tool to assist people in understanding activities that impact their carbon footprint and evaluating realistic reduction choices across energy, transport, food, waste, and water. Deliverable 1 established that urban Indian students, young professionals, and multi-generational households face acute structural barriers: split-incentive tenancy, unmetered communal water/electricity bills, extreme seasonal heat, and strict safety constraints (e.g., night-time commutes for shift workers) \cite{lee2021}. 

Mainstream carbon calculators typically calculate raw emissions using non-transparent western conversion factors, issue moralistic scoring, or suggest impractical lifestyle changes \cite{brynjarsdottir2012}. Grounded in our D1 thematic synthesis (semantic themes S1--S3 and latent themes L1--L3), our guiding Job to Be Done (JTBD) is:
\begin{quote}
\itshape
When I want to understand how my everyday routines and household habits affect the environment, I want to see which activities matter most in terms I can picture and which realistic alternatives are open to me, so that I can make changes that fit my budget, safety, rota, and shared home without being misled, lectured, or blamed for constraints I cannot control.
\end{quote}

\subsection{Core Design Principles}
To preserve rigorous traceability to our D1 empirical baseline, the system adheres to five fundamental design principles:
\begin{enumerate}[label=\textbf{DP\arabic*.},leftmargin=*]
  \item \textbf{Transparent Provenance (S1, E2):} Every numeric emission estimate must explicitly expose its formula, activity quantity, conversion factor, reference year, geographic applicability, and evidence boundary. Where no verified India-specific factor exists, the system must show an \emph{estimate unavailable} state rather than fabricating a number.
  \item \textbf{Constraint-Aware Protection (L3, E3):} Necessary routines driven by personal safety, night curfews, medical needs, or extreme weather must be flaggable as \emph{protected}, immediately disabling punitive advice or moralistic judgment.
  \item \textbf{Low-Burden, Two-Way Navigation (S3, E5):} Core estimation must require $\le 3$ primary decisions ($\le 60$ seconds). Users must be able to explore assumptions or backtrack without destructive state loss.
  \item \textbf{Communal Boundary Integrity (L1/L2, E7):} Shared utilities (e.g., central geysers, air conditioning, delivery packaging) must be displayed as collective household totals and never artificially partitioned into an individual's personal score without verified sub-metering.
  \item \textbf{Descriptive Representation for Circular Habits (S2, E4):} Indigenous Indian conservation practices (e.g., ragpickers/kabadiwala recycling, bucket baths, appliance repair) must be acknowledged descriptively without claiming unverified carbon offsets.
\end{enumerate}

\section{Ideation Process and Response Families}

To prevent premature fixation on a standard dashboard interface, the team executed a structured divergent ideation workshop utilizing the \emph{Crazy Eights} and \emph{10-plus-10} design techniques. Ideation focused on addressing the tensions identified in D1 personas: Persona 1 (Aanya, university student in a shared flat), Persona 2 (Rohan, night-shift IT professional with safety constraints), and Persona 3 (Meera, manager of a multi-generational household).

The divergent phase produced 32 raw design responses, which were affinity-clustered into eight coherent \textbf{response families}:
\begin{enumerate}[label=(\arabic*),leftmargin=*]
  \item \textbf{Contextual Explain-and-Compare Journey:} Micro-interaction connecting single activities to transparent assumptions and realistic substitutions.
  \item \textbf{Household Commons Board:} Shared living room display for coordinating communal resource use without interpersonal blame.
  \item \textbf{Quiet-Week Action Cards:} Minimalist weekly decision prompts tailored to time constraints and shift schedules.
  \item \textbf{Local Travel and Safety Explorer:} Multi-modal transport comparator prioritizing female safety and weather realities over raw carbon.
  \item \textbf{Delivery and Packaging Friction Reducer:} Tracker assessing hyper-local e-commerce packaging routines.
  \item \textbf{Assumptions and Factor Audit Explorer:} Deep-dive analytical view for evidence-conscious users verifying data sources.
  \item \textbf{Rota-Aware Appliance Planner:} Scheduler shifting heavy appliance use (washing machine, geyser) around peak grid carbon intensity.
  \item \textbf{Offline Household Conversation Cards:} Physical/card-based discussion prompts for joint decision-making between generations.
\end{enumerate}

From these eight families, the team synthesized three distinct concepts representing fundamentally different primary interaction models for detailed analysis.

\section{Three Distinct Conceptual Designs}

\subsection{Concept A: Contextual Choice Coach (Footprint Lens)}
\textbf{Interaction Paradigm:} \emph{Micro-decision checkpoint with contextual transparency.}

\begin{figure}[ht]
\centering
\begin{verbatim}
+------------------+     +-------------------+     +-------------------------+
| P0: Entry        |     | P1: Quantity      |     | P2: Contextual Result   |
| Activity, Scope, | --> | Numerical Input,  | --> | Estimate / Unavailable, |
| Protected Flag   |     | Automatic Unit    |     | Familiar Comparison     |
+------------------+     +-------------------+     +-------------------------+
                                                                |
                               +--------------------------------+-------------------------------+
                               |                                |                               |
                               v                                v                               v
                   +-----------------------+        +-----------------------+       +-----------------------+
                   | P3: Assumptions       |        | P4: Action Plan       |       | P6: Shared Boundary   |
                   | Factor, Year, Region, |        | Choose / Defer /      |       | Household total !=    |
                   | Formula, Provenance   |        | Dismiss (Constraint)  |       | Personal attribution  |
                   +-----------------------+        +-----------------------+       +-----------------------+
\end{verbatim}
\caption{Conceptual Architecture of Concept A (Contextual Choice Coach)}
\end{figure}

\textbf{Mechanics:} The user logs a single routine (e.g., electricity, water, commuting). The system displays a dual-output card pairing an illustrative carbon value with a familiar, non-carbon contextual metric (e.g., INR electricity cost or 15-litre bucket counts). Crucially, the result panel serves as a trust checkpoint: users can inspect full data provenance (P3), review an eligible next step tailored to their constraints (P4), or examine communal household boundaries (P6). Unsupported categories remain visible as descriptive records without fabricated numerical scores.

\textbf{Strengths \& Risks:} Directly resolves D1's tension between abstract carbon numbers and user action. Highly feasible in desktop JavaFX. The primary risk is false precision if illustrative factors are misinterpreted; mitigated by explicit provenance labelling.

\subsection{Concept B: Household Commons Board}
\textbf{Interaction Paradigm:} \emph{Communal coordination and conversational consensus.}

\begin{figure}[ht]
\centering
\begin{verbatim}
+-----------------------------------------------------------------------------------------------+
| HOUSEHOLD COMMONS DISPLAY (Shared Living Area / Tablet Hub)                                   |
|                                                                                               |
|  [Communal Resource Cards]                  [Weekly Household Target]                         |
|  * Home Geyser: 142 kWh (Shared)            "Move 1 AC hour/night to ceiling fan"             |
|  * Water Tanker: 3,200 Litres (Shared)      Status: In Discussion                             |
|                                                                                               |
|  [Dialogue & Action Prompts]                [Optional Consented Personal View]                |
|  "Discuss: Shift water heating to solar"    Pin-protected individual view;                    |
|  Actions: [Agree Plan] [Defer] [Table]      explicitly requires personal consent.             |
+-----------------------------------------------------------------------------------------------+
\end{verbatim}
\caption{Conceptual Layout of Concept B (Household Commons Board)}
\end{figure}

\textbf{Mechanics:} Designed for a shared tablet in communal living spaces. It aggregates whole-home utility data into shared resource cards. Instead of individual scorecards, it presents joint action proposals for household discussion. Individual views are strictly opt-in and pin-protected to prevent flatmates or parents from policing each other.

\textbf{Strengths \& Risks:} Directly addresses flatshare friction (L1) and multi-generational dynamics (L2). However, establishing verified consent boundaries in software is complex, and physical ambient displays impose high hardware constraints for an academic MVP.

\subsection{Concept C: Quiet-Week Action Cards}
\textbf{Interaction Paradigm:} \emph{Constraint-first asynchronous nudges.}

\begin{figure}[ht]
\centering
\begin{verbatim}
+-----------------------+     +-----------------------------------------------------------------+
| Constraint Filter     |     | Weekly Card Stack (Maximum 2 Cards per Week)                    |
| [x] Night Shift Rota  | --> |                                                                 |
| [x] Shared Kitchen    |     | Card 1: Shift geyser pre-heat by 30 mins before peak tariff     |
| [ ] Solo Commute      |     | Reason: Matches your evening shift rota.                        |
+-----------------------+     | Actions: [Accept Routine]   [Defer to Next Week]   [Dismiss]    |
                              +-----------------------------------------------------------------+
\end{verbatim}
\caption{Conceptual Flow of Concept C (Quiet-Week Action Cards)}
\end{figure}

\textbf{Mechanics:} The user inputs their external operational constraints (e.g., work shifts, heatwave alerts, exam periods). The system suppresses continuous daily logging and delivers a maximum of two inspectable, rule-based recommendation cards per week. Cards explain why they were selected and allow users to accept, postpone, or dismiss without penalty.

\textbf{Strengths \& Risks:} Minimizes cognitive fatigue (S3). However, recommendation logic risks appearing arbitrary or over-personalized, and without activity logging, users cannot evaluate the quantitative scale of their routines.

\section{Concept Selection and Novelty Analysis}

\subsection{Multi-Criteria Evaluation Matrix}
To determine the single concept for prototype implementation, the three models were evaluated across five weighted engineering and stakeholder criteria derived from the rubric:

\begin{table}[ht]
\centering
\small
\caption{Multi-Criteria Decision Analysis for Concept Selection}
\begin{tabularx}{\textwidth}{@{}lcp{4.2cm}p{4.2cm}p{4.2cm}@{}}
\toprule
\textbf{Criterion} & \textbf{Weight} & \textbf{Concept A: Choice Coach} & \textbf{Concept B: Commons Board} & \textbf{Concept C: Quiet Cards}\\
\midrule
D1 Must Epic Coverage & 25\% & \textbf{5/5:} Fully covers E1--E5 in a single walkthrough. & \textbf{3/5:} Weak on individual action (E3) and brief logging (E5). & \textbf{3/5:} Weak on assumptions (E2) and category coverage (E4).\\
Stakeholder Fit (P1--P3) & 25\% & \textbf{4.5/5:} Balances student flats, shift work, and family budgeting. & \textbf{4/5:} Strong for families (P3), poor for privacy-conscious students (P1). & \textbf{3.5/5:} Good for shift workers (P2), limited for general home tracking.\\
Technical Feasibility & 20\% & \textbf{5/5:} Deterministic rule engine; zero exotic hardware; JavaFX ready. & \textbf{2.5/5:} Multi-user sync, shared device authentication complex. & \textbf{3.5/5:} Requires notification engine and scheduler state.\\
Scientific Provenance & 15\% & \textbf{5/5:} Dedicated P3 assumptions table; transparent factor metadata. & \textbf{3/5:} Aggregated totals obscure underlying conversion formulas. & \textbf{2.5/5:} Recommendation rules detach action from visible evidence.\\
Risk Profile & 15\% & \textbf{4.5/5:} Low risk; false precision mitigated by explicit disclaimers. & \textbf{2/5:} High privacy/coercion risk in multi-tenant environments. & \textbf{3/5:} Risk of appearing like black-box machine learning.\\
\midrule
\textbf{Weighted Score} & \textbf{100\%} & \textbf{4.65 / 5.00 (Selected)} & \textbf{3.00 / 5.00} & \textbf{3.20 / 5.00}\\
\bottomrule
\end{tabularx}
\end{table}

\textbf{Justification of Selection:} Concept A was selected because it guarantees complete coverage of the foundational Must-have requirements (E1--E5) in a clean, highly feasible desktop architecture. It solves the primary usability barriers identified in D1—uncertainty regarding carbon numbers and guilt over unavoidable routines—while remaining fully deliverable within COMP2300 constraints.

\subsection{Novelty and Competitor Comparison}
The COMP2300 brief requires that the prototype be \emph{novel} and not substantially the same as an existing commercial product. Table~\ref{tab:competitor} benchmarks Footprint Lens against leading industry tools.

\begin{table}[ht]
\centering
\small
\caption{Competitor Analysis and Novelty Differentiation Matrix}
\label{tab:competitor}
\begin{tabularx}{\textwidth}{@{}lXXXX@{}}
\toprule
\textbf{Feature / Dimension} & \textbf{Earth Hero / Klima} & \textbf{UN Carbon Calculator} & \textbf{Capture / Giki Zero} & \textbf{Footprint Lens (Concept A)}\\
\midrule
Emission Factor Transparency & Hidden backend factors; presents single absolute kg value. & National averages; formulas and boundaries concealed. & Proprietary algorithms; black-box scoring. & \textbf{Inspectable Provenance:} Exposes factor, source, year, geography, and formula on P3.\\
Constraint / Safety Awareness & Unconditional reduction tips; penalizes car use uniformly. & Static questionnaire; ignores shift rotas or curfews. & Gamified badges; penalizes unavoidable travel. & \textbf{Protected Routine Flag:} Safety, heat, and shift constraints disable punitive advice.\\
Communal vs Personal Scope & Divides bills equally; false personal allocation. & Asks for household size and divides linearly. & Assumes individual control over all utility bills. & \textbf{Communal Boundary Integrity:} Household totals are isolated from individual shares.\\
Circular Habits \& Waste & Fabricates offset credits or ignores repair/reuse. & Only tracks flights, diet, and vehicle mileage. & Requires automated bank account transaction linking. & \textbf{Descriptive Logging:} Circular habits recorded descriptively without fake credits.\\
Algorithmic Transparency & Proprietary AI recommendations. & Static generic tips. & Black-box behavioral nudges. & \textbf{Transparent Rule Engine:} 100\% deterministic, inspectable decision logic.\\
\bottomrule
\end{tabularx}
\end{table}

\section{Small User Stories and MoSCoW Prioritisation}

To define key functionality with rigorous engineering precision, user stories are mapped to D1 Epics (E1--E5), explicit user personas (Persona 1: Aanya, flatshare student; Persona 2: Rohan, night-shift analyst; Persona 3: Meera, family manager), and prioritized using the full MoSCoW spectrum with testable acceptance criteria.

\begin{table}[ht]
\centering
\scriptsize
\caption{Comprehensive Small User Story Register with Acceptance Criteria}
\begin{tabularx}{\textwidth}{@{}p{1.2cm}p{1.3cm}p{1.5cm}Xp{4.8cm}@{}}
\toprule
\textbf{ID / Epic} & \textbf{MoSCoW} & \textbf{Stakeholder} & \textbf{User Story} & \textbf{Verifiable Acceptance Criteria}\\
\midrule
D2-US01 (E1) & \textbf{Must} & P1 (Aanya), P3 (Meera) & As a resident, I want to record an activity with its category, quantity, and scope, so that I can see what is being tracked. & \textbf{Given} an activity category, \textbf{When} a non-negative quantity is entered, \textbf{Then} unit updates automatically and scope is stored without losing form state.\\
D2-US02 (E1) & \textbf{Must} & P1--P3 & As a user, I want numeric estimates paired with familiar context (INR cost or bucket count), so that I can grasp the scale. & \textbf{Given} a valid quantity, \textbf{When} result is rendered, \textbf{Then} display kg CO2 alongside labelled familiar metric; unavailable factors never show 0.\\
D2-US03 (E2) & \textbf{Must} & P2 (Rohan), Evidence-user & As an analytical user, I want to inspect factor, formula, year, and source, so that I can verify assumptions. & \textbf{Given} any result, \textbf{When} user opens P3, \textbf{Then} table shows activity, quantity, factor label, reference year, geography, and source status.\\
D2-US04 (E2) & \textbf{Must} & P1--P3 & As a user, I want invalid inputs rejected with recovery controls, so that errors do not crash the session. & \textbf{Given} negative, blank, or $>10^6$ input, \textbf{When} submitted, \textbf{Then} specific error message displays, focus is set, and previous inputs remain intact.\\
D2-US05 (E3) & \textbf{Must} & P2 (Rohan), Shift-worker & As a shift worker, I want to mark constrained routines as protected, so that I am not blamed for necessary choices. & \textbf{Given} protected flag is checked, \textbf{When} P4 action is viewed, \textbf{Then} action buttons are disabled and explicit non-punitive explanation displays.\\
D2-US06 (E3) & \textbf{Must} & P1, P3 & As a user, I want feasible, non-preachy action options, so that I can act within my realistic boundaries. & \textbf{Given} an unprotected result, \textbf{When} action is opened, \textbf{Then} show one realistic next step with functional Choose, Defer, and Not Feasible buttons.\\
D2-US07 (E4) & \textbf{Must} & P1, P3 & As an urban resident, I want packaging, waste, and repair visible, so that circular habits are acknowledged. & \textbf{Given} packaging or repair category, \textbf{When} previewed, \textbf{Then} system displays descriptive preview (P5) without inventing carbon credits.\\
D2-US08 (E5) & \textbf{Must} & P1--P3 & As a busy user, I want a fast entry path ($\le 60$s, $\le 3$ clicks), so that logging does not waste time. & \textbf{Given} P0, \textbf{When} user loads example or enters quantity and hits Enter, \textbf{Then} result renders in $\le 2$ transitions without mandatory extra forms.\\
D2-US09 (E5) & \textbf{Must} & P1--P3 & As an assistive tech user, I want full keyboard and screen-reader support, so that I can navigate accessibly. & \textbf{Given} keyboard-only navigation, \textbf{When} navigating P0--P6, \textbf{Then} all controls have visible focus rings, status changes are announced via ARIA.\\
\midrule
D2-US10 (E1) & \textbf{Should} & P1, P3 & As a household manager, I want to view a monthly summary of logged routines, so that I can see recurring patterns. & \textbf{Given} multiple logged entries, \textbf{When} summary is opened, \textbf{Then} activities are aggregated by category and unit without cross-unit math errors.\\
D2-US11 (E5) & \textbf{Should} & P2 (Rohan) & As a privacy-conscious user, I want local CSV export of my data, so that my personal routines remain offline. & \textbf{Given} session records, \textbf{When} export is selected, \textbf{Then} standard CSV file downloads locally with zero network transmission.\\
D2-US12 (E2) & \textbf{Should} & P3 (Meera) & As a resident, I want to adjust electricity tariff rates to my state (e.g. Delhi vs Mumbai), for accurate cost comparisons. & \textbf{Given} electricity result, \textbf{When} tariff rate is modified, \textbf{Then} variable cost calculation updates immediately with an illustrative note.\\
\midrule
D2-US13 (E4) & \textbf{Could} & P1 (Aanya) & As a student, I want seasonal monsoon/heatwave routine presets, so that my logging reflects extreme weather. & \textbf{Given} seasonal preset toggle, \textbf{When} activated, \textbf{Then} water/cooling baseline recommendations adjust parameters transparently.\\
D2-US14 (E3) & \textbf{Could} & P3 (Meera) & As a consumer, I want to compare bulk vs single-use delivery packaging, to evaluate e-commerce impact. & \textbf{Given} packaging entry, \textbf{When} compared, \textbf{Then} displays descriptive packaging reduction options without carbon claims.\\
\midrule
D2-US15 (E3) & \textbf{Won't} & P1--P3 & Out of scope: Automated black-box AI predictions that infer lifestyle patterns without user consent. & \textbf{Boundary rationale:} Violates COMP2300 prohibition against speculative machine learning; compromises transparent provenance.\\
D2-US16 (E5) & \textbf{Won't} & P1, P2 & Out of scope: IoT hardware sensors or smart meter API integration. & \textbf{Boundary rationale:} Unfeasible for urban Indian rental flatshares; introduces complex third-party network dependencies.\\
D2-US17 (E3) & \textbf{Won't} & P1--P3 & Out of scope: Social leaderboards or peer-shaming rating mechanisms. & \textbf{Boundary rationale:} Directly contradicts JTBD; promotes guilt and ignores structural socio-economic constraints.\\
\bottomrule
\end{tabularx}
\end{table}

\section{Low-Fidelity Prototype Architecture and Wireframe Walkthrough}

\subsection{Fidelity, Tooling and Deliberate Boundaries}
In strict adherence to the D2 brief (``\emph{You are not expected to do any programming, and you will lose marks if the fidelity of your prototype is inappropriate... it should not have obvious conceptual or usability problems}''), the prototype is implemented as an accessible, browser-rendered interactive wireframe using pure semantic HTML5, clean CSS custom properties, and client-side JavaScript. It runs locally without external dependencies, web frameworks, or network calls.

The visual style deliberately avoids high-fidelity graphics, saturated palettes, or decorative imagery, utilizing a structured monochrome/green wireframe aesthetic that emphasizes decision paths, boundary conditions, and accessibility controls.

\subsection{Detailed Screen Walkthrough (P0--P6)}

\begin{figure}[ht]
\centering
\small
\begin{verbatim}
+-----------------------------------------------------------------------------------------------+
| FOOTPRINT LENS: D2 Working Prototype                               [Illustrative Design Only] |
+-----------------------------------------------------------------------------------------------+
| Journey Nav:    | STATE P0: START A RECORD                                                    |
| [*] P0 Start    | What would you like to understand?                                          |
| [ ] P1 Describe | Activity Category: [ Home electricity              |v]                      |
| [ ] P2 Result   | Who does this describe?: [ Me / my activity        |v]                      |
| [ ] P3 Assump.  |                                                                             |
| [ ] P4 Action   | [x] Mark as protected or necessary (safety, curfew, shift work, heat)       |
| [ ] P5 Unavail. | Note: This category has an illustrative factor only; inspect assumptions.  |
| [ ] P6 Shared   |                                                                             |
|                 | [ Continue to quantity ]   [ Load illustrative example ]   [ Shared preview ] |
| [Reset Prototype]                                                                             |
+-----------------------------------------------------------------------------------------------+
\end{verbatim}
\caption{Wireframe Layout for Panel P0 (Activity Selection, Scope and Constraint Flagging)}
\end{figure}

\begin{itemize}[leftmargin=*]
  \item \textbf{P0: Start a Record (Entry):} The user selects one of seven D1-relevant categories (\emph{electricity, water, travel, food, packaging, waste, repair}), selects personal or household scope, and can flag the routine as \emph{protected}. A pre-loaded illustrative example (168 kWh household electricity) enables reproducible assessment walkthroughs.
  
  \item \textbf{P1: Describe It (Quantity Input):} Accepts numerical quantity. The unit is strictly owned and displayed by the category (e.g., kWh, litres, passenger-km, kg) via a dedicated, styled read-only \code{<output>} element to prevent incompatible calculations. Input supports decimal values and the standard keyboard \code{Enter} key. Missing, negative, or $>10^6$ values trigger immediate, non-destructive validation announcements.
  
  \item \textbf{P2: See the Result (Universal Result State):} The universal evaluation checkpoint. Displays the calculated estimate ($120.288\text{ kg CO}_2$) or an explicit ``Estimate unavailable'' notification. Crucially, it renders a familiar contextual equivalence (e.g. \emph{Variable cost context: \rupee 1,344.00}) and links to P3 (Assumptions), P4 (Realistic Action), P6 (Shared Total), and P1 (Edit Quantity).
  
  \item \textbf{P3: Inspect Assumptions (Provenance Transparency):} Renders an evidence table detailing Activity Name, User Quantity, Stored Scope, Factor Value ($0.716\text{ kg CO}_2\text{/kWh}$), Factor Source (D1 fixture), Reference Year (2026 illustrative), Geography (India), and Calculation Formula. Explicitly highlights evidence boundaries.
  
  \item \textbf{P4: Choose an Action (Constraint-Aware Decision):} Displays a realistic, non-preachy next step. When the routine was flagged as \emph{protected} on P0, the \emph{Choose}, \emph{Defer}, and \emph{Not Feasible} buttons are gracefully disabled, accompanied by an explicit status notice: ``\emph{Action decisions are disabled because this activity is marked protected or necessary.}''
  
  \item \textbf{P5: Unknown Estimate (Category Boundary):} Dedicated recovery route for categories lacking verified conversion factors (e.g., delivery packaging, circular repair). Explains what evidence is missing and affirms that keeping a descriptive record is valuable without fabricating false carbon offsets.
  
  \item \textbf{P6: Shared Total Boundary (Communal Protection):} Features a visual boundary diagram demonstrating that a communal household total does not equal or imply an individual's personal share. Offers a non-destructive ``Back to result'' route.
\end{itemize}

\subsection{Accessibility, Technical Feasibility and Verification}
The prototype incorporates comprehensive accessibility engineering:
\begin{itemize}[leftmargin=*]
  \item \textbf{Keyboard Operability:} Full tab-stop sequence with distinct 3px focus outlines (\code{:focus-visible}); form submission executable via the \code{Enter} key.
  \item \textbf{Screen Reader Compatibility:} Dynamic status changes and errors announced via an \code{aria-live="assertive"} polite region; input errors mapped with \code{aria-describedby} and \code{aria-invalid}. Panel headings dynamically receive focus (\code{tabindex="-1"}) upon route changes.
  \item \textbf{Visual Clarity:} Non-reliance on color alone; high-contrast ratios exceeding WCAG 2.1 AA ($>4.5:1$ for text); responsive single-column layout wrapping on viewports down to 320px width.
\end{itemize}
Technical feasibility is validated through automated test coverage: eight pure logic unit tests (\code{prototype-logic.test.mjs}) verifying unit compatibility, calculations, protected states, and boundaries; a structural contract check (\code{prototype-structure-check.py}) confirming all IDs, panels, and handlers; and clean local HTTP serving.

\section{User Evaluation Protocol for D3/D7}

To satisfy the rubric requirement for a robust evaluation plan, we specify a controlled, task-based usability study designed for execution in Deliverable 3:

\subsection{Participant Selection and Ethics}
Testing will involve $N=6$ participants recruited from peers across other COMP2300 project groups, matching our target demographics (student flatshare tenants, working professionals, and multi-generational householders). To ensure ethical integrity, no sensitive personal energy bills, financial records, or location histories will be collected; all tasks will utilize standardized synthetic scenarios. Participants will complete an informed consent form and debriefing.

\subsection{Structured Task Scenarios}
\begin{enumerate}[label=\textbf{Task \arabic*:},leftmargin=*]
  \item \textbf{Standard Activity Logging and Provenance Check (E1, E2):} Participant is given a synthetic utility bill showing 150 kWh electricity. They must log the entry as personal scope, navigate to P2, open P3, and correctly identify the conversion factor and data year.
  \item \textbf{Constraint Protection (E3):} Participant is asked to log a late-night auto-rickshaw commute (15 passenger-km) undertaken for personal safety. They must mark the routine as protected, proceed to P4, and explain why action buttons are disabled.
  \item \textbf{Descriptive Logging for Unsupported Category (E4):} Participant logs a routine delivery packaging entry (3 kg). They must navigate through the result, open P5, and verify that the system does not fabricate a numerical carbon credit.
  \item \textbf{Shared Boundary Comprehension (E1, E7):} Participant inspects a 500 kWh household electricity total, navigates to P6, and answers an explain-back comprehension question regarding whether this total can be attributed to them individually.
\end{enumerate}

\subsection{Evaluation Metrics}
\begin{itemize}[leftmargin=*]
  \item \textbf{Effectiveness:} Task Completion Rate (binary success/failure per task); Critical Error Rate (submitting invalid data without recovery).
  \item \textbf{Efficiency:} Time on Task (target: $\le 60$ seconds for Task 1); Total interaction step count.
  \item \textbf{Cognitive Comprehension:} Post-task explain-back probes verifying whether participants understood that illustrative numbers are not verified emissions.
  \item \textbf{Usability and Satisfaction:} Single Ease Question (SEQ, 1--7 scale) after each task; System Usability Scale (SUS) questionnaire administered post-session, with a target benchmark score of $\ge 75/100$.
\end{itemize}

\section{Project Management, GitLab Workflow and Process Reflection}

\subsection{Team Working Methods and Collaboration}
In accordance with our Week 6 and Week 7 minutes, team responsibilities were distributed across complementary specializations:
\begin{itemize}[leftmargin=*]
  \item \textbf{Arnav Mishra (Research, Requirements \& Report):} Coordinated D1-to-D2 traceability, authored formal LaTeX documentation, formulated user stories, and synthesized the evaluation protocol.
  \item \textbf{Keshav Singhal (Engineering, Architecture \& Logic):} Implemented the deterministic business logic engine, automated unit tests, MoSCoW prioritization matrix, and technical feasibility verification.
  \item \textbf{Hiranya Makkar (Interaction Design, Prototype \& Accessibility):} Designed wireframe layouts, implemented accessible HTML/CSS/JS prototype panels, and audited WCAG 2.1 AA keyboard/screen-reader compliance.
\end{itemize}
Collaboration was maintained through a shared GitLab repository, structured branch workflows (\code{feature/d2-conceptual-design}), issue tracking, and weekly peer review rotations.

\subsection{Critical Process Reflection}
The transition from D1 to D2 underscored the vital distinction between user intentions and structural feasibility. During initial ideation, the team considered incorporating automated electricity bill scanning and ambient shared displays (Concept B). However, critical evaluation revealed that bill scanning introduces brittle OCR dependencies and privacy hurdles for flatmates, while ambient boards require speculative hardware. Selecting Concept A allowed the team to deliver an honest, rigorous design artefact that confronts the core psychological challenges of sustainability—transparent data and constraint awareness—without technological over-promising.

\section{Conclusion and Deliverable Release Gate}

Deliverable 2 successfully bridges the empathic baseline of D1 into a concrete, testable conceptual architecture. By formulating three distinct concepts, selecting Concept A through an objective decision matrix, specifying comprehensive user stories across the full MoSCoW spectrum, and engineering a functional, accessible low-fidelity prototype, the project establishes a rigorous foundation for implementation in D3 and D4. All components have been verified through automated unit testing, structural contract checks, and formatting audits, fulfilling all criteria for Distinction-level engineering practice.

\begin{thebibliography}{9}
\bibitem{hoffmann2024} S. Hoffmann, W. Lasarov, H. Reimers and M. Trabandt. 2024. Carbon footprint tracking apps: What drives user engagement and behavioral change? \emph{Journal of Cleaner Production} 434, 139981. DOI: \href{https://doi.org/10.1016/j.jclepro.2023.139981}{10.1016/j.jclepro.2023.139981}.
\bibitem{lasarov2024} W. Lasarov, S. Hoffmann, R. Mai and J. Schleich. 2024. Carbon footprint tracking apps: the spillover effects of feedback and goal-activating appeals. \emph{Ecological Economics} 226, 108354. DOI: \href{https://doi.org/10.1016/j.ecolecon.2024.108354}{10.1016/j.ecolecon.2024.108354}.
\bibitem{brynjarsdottir2012} H. Brynjarsdottir, M. Hakansson, J. Pierce, E. Baumer, C. DiSalvo and P. Sengers. 2012. Sustainably unpersuaded: How persuasion narrows our vision of sustainability. In \emph{Proc. ACM CHI '12}. ACM, 947--956. DOI: \href{https://doi.org/10.1145/2207676.2208539}{10.1145/2207676.2208539}.
\bibitem{lee2021} J. Lee, O. Taherzadeh and K. Kanemoto. 2021. The scale and drivers of carbon footprints in households, cities and regions across India. \emph{Global Environmental Change} 66, 102205. DOI: \href{https://doi.org/10.1016/j.gloenvcha.2020.102205}{10.1016/j.gloenvcha.2020.102205}.
\bibitem{mohanty2023} V. Mohanty, A. Filipowicz, N. Bravo, S. Carter and D. A. Shamma. 2023. Save a tree or 6 kg of CO2? In \emph{Proc. ACM CHI '23}. ACM. DOI: \href{https://doi.org/10.1145/3544548.3580675}{10.1145/3544548.3580675}.
\bibitem{shi2021} S. Shi and J. Yin. 2021. Global research on carbon footprint: A scientometric review. \emph{Environmental Impact Assessment Review} 89, 106571. DOI: \href{https://doi.org/10.1016/j.eiar.2021.106571}{10.1016/j.eiar.2021.106571}.
\end{thebibliography}

\end{document}
